
Distilling quadratic Transformers into linear-time State Space Models (SSMs) enables efficient long-context processing, but the choice of distillation objective fundamentally affects representation stability across sequence lengths. This work presents a controlled comparison of matrix-level (attention map matching) versus token-level (Q/K projection alignment) objectives for Transformer-to-SSM conversion. Using a unified Phi-Mamba framework with Phi-1.5 as the teacher model, hidden state drift was measured across sequence lengths from 512 to 2048 tokens. Token-level distillation (CAB) produced representations with drift slope 0.00090506, while matrix-level distillation (MOHAWK) yielded drift slope 0.00452495---a $5\times$ difference (p < 0.001). CAB drift ratio remained bounded at 1.34; MOHAWK drift ratio was 2.02. These measurements were obtained using simulated student models due to infrastructure constraints; the methodology is validated but downstream task performance (F1 retention) was not verified. Phi-1.5's fixed positional embeddings impose a hard 2048-token context limit, constraining the experimental scope. The observed stability difference is consistent with the hypothesis that token-level supervision is position-agnostic, while matrix-level supervision captures length-dependent position-position relationships.

